% Desenvolvimento
\chapter{EXPOSIÇÃO DO TEMA OU MATÉRIA}

É a parte principal e mais extensa do trabalho.
Deve apresentar a fundamentação teórica,
a metodologia,
os resultados
e a discussão.
Divide-se em seções e subseções
conforme a NBR 6024~\cite{abnt14724}.
Quanto a sua estrutura,
segue as recomendações da norma
para preparação de trabalhos acadêmicos,
a NBR 14724 de 2011~\cite{abnt14724}.
Quanto à Formatação,
segue o modelo adotado pela UFSC, o formato A5.

\section{FORMATAÇÃO DO TEXTO}

No que diz respeito à estrutura do trabalho,
o novo modelo para dissertações e teses
adotado pela UFSC
segue a NBR 14724 (2011).
Porém, em relação à formatação,
a UFSC adotou o tamanho A5,
que corresponde à metade do A4.
Por esta razão,
foi necessário uma adequação no tamanho da fonte,
espaçamento entrelinhas,
margens,
etc,
conforme exposto no quadro abaixo.

O texto deve ser justificado,
digitado em cor preta,
podendo utilizar outras cores
somente para as ilustrações.
Utilizar papel branco.
Os elementos pré-textuais
devem iniciar no anverso da folha,
com exceção da ficha catalográfica.
Os elementos textuais e pós-textuais
devem ser digitados no anverso e verso das folhas,
com espaçamento simples (1).

\subsection{As ilustrações}

Independente do tipo de ilustração 
(quadro, desenho, figura, fotografia, mapa, entre outros)
sua identificação aparece na parte superior,
precedida da palavra designativa.

A indicação da fonte consultada
deve aparecer na parte inferior,
elemento obrigatório mesmo que seja produção do próprio autor.
A ilustração deve ser citada no texto
e inserida o mais próximo possível do texto
a que se refere~\cite{abnt14724}.

A Figura \ref{fig:a} mostra um quadrado.
\begin{figure}[!htb]
   \centering
   \caption{Quadrado.}\label{fig:a}
   \begin{tikzpicture}
        \draw (0, 0) rectangle (1, 1);
    \end{tikzpicture}
\end{figure}

A Tabela~\ref{tab:a} mostra mais informações do template BU.

\begin{table}[!htb]
\begin{center}
    \caption{Formatação do texto.}\label{tab:a}
    \begin{tabular}{ p{3cm} | p{6cm} }
        \hline
        Cor & Branco \\
        \hline
        Formato do papel & A5 \\
        \hline
        Gramatura & 75 \\
        \hline
        Impressão & Frente e verso \\
        \hline
        Margens & Espelhadas: superior 2, Inferior: 1,5, Externa 1,5 e Externa: 2. \\
        \hline
        Cabeçalho & 0,7 \\
        \hline
        Rodapé & 0,7 \\
        \hline
        Paginação & Externa \\
        \hline
        Alinhamento vertical & Superior \\
        \hline
        Alinhamento do texto & Justificado \\
        \hline
        Fonte sugerida & Times New Roman  \\
        \hline
        Tamanho da fonte & 10,5 para o texto
        incluindo os títulos das seções e subseções.
        As citações com mais de três linhas
        as legendas das ilustrações e tabelas, fonte 9,5. \\
        \hline
        Espaçamento entre linhas & Um (1) simples \\
        \hline
        Espaçamento entre parágrafos & Anterior 0,0; Posterior 0,0 \\
        \hline
        Numeração da seção & As seções primárias devem começar
        sempre em páginas ímpares.
        Deixar um espaço (simples) entre o título da seção
        e o texto e entre o texto e o título da subseção. \\
        \hline
    \end{tabular}
\end{center}
Fonte: Universidade Federal de Santa Catarina (2011)
\end{table}

\subsection{Equações e fórmulas}

As equações e fórmulas
devem ser destacadas no texto
para facilitar a leitura.
Para numerá-las,
deve-se usar algarismos arábicos entre parênteses
e alinhados à direita.
Pode-se usar uma entrelinha maior do que a usada no texto~\cite{abnt14724}.

Exemplo: A equação \ref{eq:a}
\begin{equation}
    x^2 + y^2 = z^2
    \label{eq:a}
\end{equation}
 e a equação \ref{eq:b}
\begin{equation}
    x^2 + y^2 = n
    \label{eq:b}
\end{equation}

\section{EXEMPLO DE COMO GERAR A LISTA DE SÍMBOLOS E ABREVIATURAS}

Para gerar a lista de símbolos e abreviaturas use os comandos

\abreviatura{ABNT}{Associação Brasileira de Normas Técnicas}
\abreviatura{IBGE}{Instituto Brasileiro de Geografia e Estatística}
\simbolo{$\int$}{Integral}
\simbolo{$\prod$}{Produtório}

\begin{lstlisting}
\simbolo{símbolo}{descrição}
\end{lstlisting}

\begin{lstlisting}
\abreviatura{abreviatura}{descrição}
\end{lstlisting}

\subsection{Exemplo de citações no \LaTeX}

Segundo \citeonline{alves_2001} ...

...no final da frase \cite{abnt14724,BU_formatoA5}

\nocite{alves_2001,abnt10520,abnt6024,abnt14724}
